\maketitle
\tableofcontents

% ============================================================
\chapter*{Dok.~313 -- Kein Anfang}
\addcontentsline{toc}{chapter}{Dok.~313 -- Kein Anfang}
% ============================================================

\section*{Worum es geht}

Dok.~312 endet mit Pflicht~D(ii): Können die Feldkonfigurationen
auf dem kompakten Zeitzyklus ($\tau_c = 2\pi/H_0 = 91{,}9$~Gyr)
periodisch schließen \emph{und dabei} die beobachtete thermische
Abfolge tragen -- Nukleosynthese, Rekombination, Strukturbildung?
Ein Kreis hat keinen Anfang; die Beobachtungen zeigen aber einen
Gradienten. Dieses Dokument rechnet die Frage durch
(\texttt{ffgft\_313\_zeitzyklus\_geschichte.py},
\texttt{ffgft\_313\_fraktale\_wegkorrektur.py}) und findet drei
Ergebnisse und einen harten offenen Kern:

\begin{enumerate}
  \item \textbf{Antipoden-Befund [K$|$P20]:} Die beobachtbare
        Geschichte füllt \emph{einen Halbzyklus}. Mit der
        fraktalen Wegkorrektur (Kap.~B, unabhängig aus A040)
        sitzt der Partikelhorizont auf $0{,}14\,\%$ am Antipoden
        ($\theta = 0{,}9986\,\pi$; glatt gerechnet
        $1{,}012\,\pi$), die Streuwand bei $0{,}979\,\pi$ --
        der \enquote{Anfang} ist der Antipode des Zeitzyklus,
        kein Ereignis.
  \item \textbf{Glatte Schließung [K]:} Das beobachtete
        Massenlauf-Profil erreicht den Antipoden mit Steigung
        $\sqrt{\Omega_r} \approx 0{,}01$ -- die
        Glattheitsbedingung einer geraden periodischen Schließung
        ist auf $1\,\%$ erfüllt, und die Rückkehr-Hälfte ist exakt
        die photonisch unbeobachtbare Hälfte.
  \item \textbf{Entropie-Hindernis [K]:} Exakte
        $\tau_c$-Periodizität ist dynamisch nicht generisch
        (Poincaré-Diskrepanz $\sim 10^{10^{104}}$); sie muss als
        Randbedingung gesetzt werden. Feinkörnig zulässig,
        grobkörnig offen -- das ist der präzisierte Kern von
        D(ii).
\end{enumerate}

Kein Teil der A-Serie; Registrierung in Dok.~190 separat. Bezüge:
Dok.~306 ($\partial\text{T}^4 = \emptyset$), 267 (Entartung),
312 (Abschlussskala, kompakte Zeitrichtung, Statik der Relation,
Torus-Ruhesystem).

% ============================================================
\chapter{Kein Rand, kein $t=0$: die Ausgangslage}
% ============================================================

Die Position der FFGFT zu Singularität und Urknall steht seit
Dok.~306: $\partial\text{T}^4 = \emptyset$ -- der Torus hat keinen
Rand, \enquote{was war bei $t=0$?} ist falsch gestellt. Dok.~312
hat diese Position konkretisiert und dreifach gestützt:

\textbf{Geometrisch:} Die Zeitrichtung ist ein Zyklus der Länge
$91{,}9$~Gyr. Ein geschlossener Kreis hat keinen Anfangspunkt --
\enquote{$t=0$} ist auf ihm so sinnlos wie \enquote{der Anfang des
Äquators}.

\textbf{Rückwärts:} In der Massenlauf-Lesart (Dok.~312, Statik der
Relation; Wetterich) gehen zurückgerechnet nicht Abstände gegen
null, sondern Massen: $\tilde T \cdot m = 1$ heißt dann
$\tilde T \to \infty$ -- keine Dichtedivergenz, sondern ein Regime
immer langsamerer Uhren. Die \enquote{Singularität} ist der Ort,
an dem die Uhren stehen, nicht an dem die Dichte divergiert.

\textbf{Vorwärts:} Kein Kollaps in eine Singularität -- das
Windungs-/Impuls-Minimum (Anhang Dok.~312) hat $V'' > 0$.

Was Dok.~312 offen ließ, ist die \emph{Geschichte}: Die
Beobachtungen zeigen eine klare thermische Abfolge. Wie verträgt
sich ein Gradient mit einem Kreis? Das ist D(ii), und es wird hier
gerechnet, nicht erzählt.

% ============================================================
\chapter{Vorschaltung: die fraktale Wegkorrektur}
% ============================================================

Alle Distanzen dieses Dokuments sind \emph{Lichtwege}. Der Raum
der FFGFT ist aber nicht glatt-affin, sondern trägt die fraktale
Dimension $D_f = 3-\xi$ (A040). Ein durchlaufener Weg akkumuliert
darüber eine Verlängerung; die glatt gerechnete Distanz ist zu
groß. Vor allen Zahlen steht deshalb die Frage, wo dieser Faktor
greift.

\section{Der Faktor}

A040 liefert ihn, und zwar \emph{abgeleitet}, nicht gesetzt:
\begin{equation}
  K_\text{frak} = \Bigl(\frac{D_f^\text{eff}}{3}\Bigr)^{D_f^\text{eff}/2}
  = 1-100\,\xi = 0{,}98665,
  \qquad D_f^\text{eff} \approx 2{,}973 .
\end{equation}
$D_f^\text{eff}$ ist der einzige Wert, für den zwei unabhängige
Herleitungen von $m_e/m_\mu$ übereinstimmen -- der Faktor $100$
ist damit eine abgeleitete Größe, \emph{keine} Dekadenzählung und
nichts, was man an ein gewünschtes Ergebnis anpassen könnte. Seine
Unsicherheit ist beziffert: $n = 100 \pm 2$ (Register~190, R67),
also $\pm 0{,}027\,\%$ im Effekt. Die Wegverlängerung beträgt
$1/K_\text{frak} - 1 = 1{,}35\,\%$.

\section{Die Zuordnungsregel}

A040 schränkt scharf ein [B]: $K_\text{frak}$ betrifft
\emph{absolute} Werte, niemals Verhältnisse \emph{derselben
Ebene} -- dort kürzt er sich heraus. Was \enquote{dieselbe Ebene}
heißt, entscheidet die Unterscheidung aus Dok.~311:

\begin{center}
\begin{tabular}{lll}
\toprule
Größe & Ebene & $K_\text{frak}$ \\
\midrule
Lichtweg $\eta_0$, $\chi(z)$ & ausgerollt (durchlaufen) & greift \\
$L^{*}$, $R_H$, $\tau_c$ & eingerollt (topologisch) & --- \\
Weg $/$ Struktur ($\eta_0/\pi R_H$) & gemischt & \textbf{greift} \\
Weg $/$ Weg ($\chi_\text{rec}/\eta_0$) & ebenengleich & kürzt \\
Winkel, dimensionslose Verhältnisse & ebenengleich & kürzt \\
\bottomrule
\end{tabular}
\end{center}

Der Grund ist begrifflich, nicht numerisch: Ein Lichtweg existiert
nur in der \emph{ausgerollten} Darstellung -- in der eingerollten
Sicht gibt es keinen \enquote{Weg zum Antipoden}, nur die
Zyklusposition $\theta$. $R_H$ dagegen folgt aus
$H_0 = (\pi/2)c\,\xi^{10}/\bar\lambda_e$, einer topologischen
Strukturrelation. $\eta_0/(\pi R_H)$ vergleicht also Durchlaufenes
mit Topologischem und ist \emph{kein} ebenengleiches Verhältnis.

\textbf{Kontrolltest:} $\chi_\text{rec}/\eta_0 = 0{,}98024$ ist
mit und ohne Korrektur identisch -- zwei Lichtwege, der Faktor
kürzt sich wie gefordert. Die Regel ist also nicht frei
anwendbar, sondern hat eine prüfbare Invariante.

\textbf{Deklaration:} Diese Zuordnung ist eine begriffliche
Setzung, und sie verbessert die Zahlen. Beides gehört
nebeneinander gesagt. Für sie spricht, dass die Wirkungsrichtung
physikalisch zwingend ist (ein gewundener Weg überbrückt bei
gleicher Laufzeit weniger glatte Distanz -- keine Wahl) und dass
die Ebenen-Unterscheidung aus Dok.~311 unabhängig und früher
formuliert wurde. Gegen sie spricht nichts Bekanntes, aber sie ist
nicht bewiesen. Alle folgenden Zahlen werden deshalb in beiden
Fassungen geführt.

% ============================================================
\chapter{Der Antipoden-Befund [K$|$P20]}
% ============================================================

\section{Wo sitzt die Geschichte auf dem Zyklus?}

Position auf dem Zyklus: $\theta = \chi/R_H$, mit $\chi$ als
Lichtweg-Distanz (die $z \to \chi$-Zuordnung übernimmt die
gemessene $\Lambda$CDM-Abbildung; nur dimensionslose Verhältnisse
gehen ein, Dok.~267). Der Antipode liegt bei $\theta = \pi$, d.\,h.
$\chi = L^{*}/2 = \pi R_H = 14{,}09$~Gpc. Gerechnet:

\begin{center}
\begin{tabular}{lccc}
\toprule
 & \multicolumn{2}{c}{Anteil von $L^{*}/2$} & \\
 & glatt & fraktal & \\
\midrule
Rekombination ($z=1090$) & $0{,}992$ & $0{,}979$ & $-2{,}1\,\%$ \\
Partikelhorizont ($z\to\infty$) & $1{,}012$ & $\mathbf{0{,}9986}$
  & $\mathbf{-0{,}14\,\%}$ \\
\bottomrule
\end{tabular}
\end{center}

\noindent Der \textbf{Partikelhorizont} -- die formale Grenze des
Kausalkontakts und damit der eigentliche \enquote{Anfang} --
landet mit der fraktalen Korrektur auf $0{,}14\,\%$ am Antipoden,
gegenüber $1{,}2\,\%$ in der glatten Rechnung. Das ist eine
Verbesserung um den Faktor~9, ohne dass etwas angepasst wurde. Der
Rest ist allerdings das Fünffache der $K_\text{frak}$-Unsicherheit
($\pm0{,}027\,\%$), also kein exakter Treffer.

Die \textbf{Streuwand} wandert dabei von $0{,}992$ auf $0{,}979$:
Der Nebenbefund \enquote{Rekombination am Antipoden} wird
\emph{schwächer}. Beide Aussagen können nicht gleichzeitig exakt
sein -- sie unterscheiden sich um den invarianten Faktor
$\chi_\text{rec}/\eta_0 = 0{,}980$.

\begin{center}
\begin{tabular}{lccc}
\toprule
$z$ & $\theta/\pi$ glatt & $\theta/\pi$ fraktal
  & $m(z)/m_0$ \\
\midrule
$0{,}5$ & $0{,}140$ & $0{,}138$ & $0{,}667$ \\
$1$ & $0{,}243$ & $0{,}240$ & $0{,}500$ \\
$3$ & $0{,}465$ & $0{,}459$ & $0{,}250$ \\
$10$ & $0{,}689$ & $0{,}680$ & $0{,}091$ \\
$100$ & $0{,}917$ & $0{,}905$ & $0{,}0099$ \\
$1090$ & $0{,}992$ & $0{,}979$ & $0{,}00092$ \\
\bottomrule
\end{tabular}
\end{center}

\textbf{Befund:} Die gesamte beobachtbare Geschichte -- bis
$z \to \infty$ -- füllt auf $1{,}2\,\%$ genau einen
\emph{Halbzyklus}. Die letzte Streufläche sitzt am Antipoden
($99{,}2\,\%$), der formale \enquote{Anfang} $1{,}2\,\%$ dahinter.
Der \enquote{Urknall} ist damit ein \emph{Ort auf dem Zyklus} --
die Antipodenregion, in der Massenlauf-Lesart: die Region der
kleinsten Massen und langsamsten Uhren -- und kein Ereignis vor
der Zeit.

\section{Ehrliche Einordnung der Koinzidenz [S]}

In $\Lambda$CDM-Sprache ist derselbe Sachverhalt die bekannte
Quasi-Koinzidenz $\eta_0 H_0/c \approx \pi$ (hier: $3{,}18$ gegen
$3{,}14$, also $1{,}2\,\%$) -- dort eine zufällige Folge der
gemessenen $\Omega$-Werte. In der kompakten Lesart wird daraus
eine Strukturaussage: \emph{Die Streuwand sitzt am Antipoden.}
Damit sie mehr ist als eine umgedeutete Koinzidenz, müsste die
Lage vorwärts folgen. Die fraktale Wegkorrektur ist ein Schritt in
genau diese Richtung: Sie ist unabhängig abgeleitet (A040), nicht
an dieses Ergebnis angepasst, und sie reduziert den Abstand des
Partikelhorizonts zum Antipoden von $1{,}2\,\%$ auf
$0{,}14\,\%$. Vollständig hergeleitet ist die Lage damit nicht --
der verbleibende Rest ist das Fünffache der R67-Unsicherheit, und
warum der Materiegehalt gerade den Halbzyklus füllt, bleibt offen.
Status [S], als Herleitungsziel notiert.

% ============================================================
\chapter{Die glatte Schließung [K]}
% ============================================================

\section{Das Profil und seine Endpunktsteigung}

Periodizität verlangt, dass der Massenlauf $m(\theta)$ auf dem
Kreis stetig schließt. Die einfachste Schließung ist ein
\emph{gerades} Profil um den Antipoden: Minimum bei $\theta=\pi$,
symmetrische Rückkehr auf der anderen Hälfte. Glatt ist sie, wenn
$dm/d\theta \to 0$ am Umkehrpunkt.

In der Strahlungsära gilt $H \simeq H_0\sqrt{\Omega_r}(1{+}z)^2$,
daraus analytisch

\[
  \eta_0 - \chi = \frac{c}{H_0\sqrt{\Omega_r}}\,\frac{1}{1{+}z}
  \qquad\Longrightarrow\qquad
  \left|\frac{dm}{d\theta}\right|_{\text{Ende}}
  = \sqrt{\Omega_r} \approx 0{,}0096.
\]

\textbf{Befund:} Das beobachtete Profil läuft mit Steigung
$\sim 0{,}01$ in den Endpunkt -- die Glattheitsbedingung einer
geraden periodischen Schließung ist auf $1\,\%$ erfüllt. Die
nötige Deformation gegenüber der gemessenen Geschichte ist ein
$1\,\%$-Effekt am äußersten Rand.

\section{Warum das kein Selbstbetrug ist}

Die Rückkehr-Hälfte -- dort, wo der Massenlauf wieder ansteigen
müsste -- ist \emph{exakt} die Hälfte hinter der Streuwand:
photonisch prinzipiell unbeobachtbar. Das ist zunächst verdächtig
bequem, aber es ist keine Wahl dieses Dokuments, sondern Folge des
Antipoden-Befunds: Die Wand \emph{liegt} bei $0{,}992\,\pi$. Die
Schließung widerspricht deshalb keiner einzigen Photon-Beobachtung
-- und sie ist trotzdem nicht unprüfbar (Kap.~E).

Das Klimazonen-Bild aus der Diskussion zu Dok.~312 wird damit
quantitativ: Der Gradient ist die Landkarte \emph{entlang} des
Zyklus (wie Klimazonen entlang eines Meridiankreises), keine
Entwicklung \emph{aus} einem Anfang -- und der \enquote{heißeste}
Punkt der Karte ist der Antipode.

% ============================================================
\chapter{Das Entropie-Hindernis: der Kern von D(ii) [K]}
% ============================================================

\section{Periodizität ist nicht generisch}

Das Entropiebudget des beobachtbaren Volumens ist
$S \sim 10^{104}\,k_B$ (dominiert von supermassereichen Schwarzen
Löchern; Photonen: $10^{89}\,k_B$). Poincaré-Wiederkehr eines
generischen Zustands braucht $\sim e^{S/k_B} \sim 10^{10^{104}}$
dynamische Zeiten -- gegen $\tau_c \sim 10^{18{,}5}$~s. Die
Diskrepanz beträgt $\sim 10^{104}$ \emph{Größenordnungen im
Exponenten}: Exakte $\tau_c$-Periodizität kann dynamisch nicht
erwartet werden.

\textbf{Konsequenz:} Periodizität ist eine
\textbf{Randbedingung} -- eine Auswahl periodischer Lösungen aus
dem Lösungsraum, keine Folge der Evolution. Das ist dieselbe
Kategorie wie die Wahl der T$^4$-Topologie selbst ([SETZUNG] im
Sinne von Dok.~311) und muss so gebucht werden.

\section{Fein- und grobkörnig getrennt}

\textbf{Feinkörnig [K]:} Unter unitärer Evolution ist die
feinkörnige Entropie konstant; einer exakt periodischen Lösung
steht mikroskopisch nichts entgegen. Die Randbedingung ist
zulässig.

\textbf{Grobkörnig [offen]:} Der Zweite Hauptsatz ist eine
Aussage über grobkörnige Entropie relativ zu einem Beobachter.
Auf der Rückkehr-Hälfte muss sie abnehmen -- was beobachterrelativ
möglich ist (Umkehr der Korrelationsrichtung; keiner
\emph{dort} sähe eine Verletzung, da der thermodynamische
Zeitpfeil mitdreht), aber unbewiesen. \textbf{Genau das ist der
präzisierte Kern von D(ii):} nicht die Feldperiodizität (die ist
als Randbedingung konsistent, Kap.~C), sondern der Nachweis, dass
eine grobkörnige Entropie-Schließung mit der beobachteten
Strukturbildung verträglich ist. Bis dahin: offen, ohne
Beschönigung.

% ============================================================
\chapter{Hawking-Strahlung auf dem Zeitzyklus}
% ============================================================

Passt die Hawking-Strahlung in dieses Bild? Ja -- und zwar besser
als zuvor, weil die kompakte Zeitrichtung (Dok.~312) ihr unbemerkt
den Unterbau geliefert hat. Zugleich verschärft sie den
Entropie-Kern von D(ii). Beides gerechnet
(\texttt{ffgft\_313\_hawking\_zyklus.py}).

\section{Ein Prinzip, drei Gesichter [K]}

Die drei berühmten \enquote{Temperaturen aus dem Nichts} --
Unruh, Gibbons--Hawking, Hawking -- sind im Lehrbuch drei separate
Wunder mit drei Herleitungen. Dahinter steht \emph{eine}
KMS-Regel: Wo die Zeitstruktur periodisch ist, liest ein
Beobachter Temperatur ab, $T = \hbar/(k_B\tau)$:

\begin{center}
\begin{tabular}{lll}
\toprule
System & Periode $\tau$ & $T = \hbar/(k_B\tau)$ \\
\midrule
Kosmos (kompakte Zeit) & $2\pi/H_0 = 2{,}9\times10^{18}$~s &
  $2{,}63\times10^{-30}$~K \\
Schwarzes Loch ($M_\odot$) & $8\pi GM/c^3 = 1{,}24\times10^{-4}$~s &
  $6{,}17\times10^{-8}$~K \\
Unruh ($a = g$) & $2\pi c/a = 1{,}9\times10^{8}$~s &
  $3{,}98\times10^{-20}$~K \\
\bottomrule
\end{tabular}
\end{center}

Dok.~312 hat diese Regel für den Kosmos \emph{geerdet}: Die
Gibbons--Hawking-Temperatur folgt aus der Kompaktheit der
Zeitrichtung selbst, ohne Horizont. Damit ist die
Hawking-Strahlung kein Fremdkörper und kein Extra-Postulat -- sie
ist der \emph{lokale Spezialfall derselben Regel}.

\section{Die Membran bekommt ihr Thermometer [S]}

Das Membranbild (Narrativ Kap.~04: das Loch \emph{ist} seine
fraktale Membran, die Strahlung kommt von der schwingenden
Oberfläche) hatte bisher keine Antwort auf \enquote{warum genau
diese Temperatur?}. Jetzt gibt es sie, in der Sprache der
Massenlandkarte: Die Masse staucht die Zeitstruktur
($\tilde T \cdot m = 1$ -- je näher der Membran, desto langsamer
die Uhren), und an der Membran wird die lokale Zeitstruktur
periodisch mit $\tau = 8\pi GM/c^3$. Kurze Periode, hohe
Temperatur:

\begin{center}
\begin{tabular}{ll}
\toprule
Objekt & $T_H$ \\
\midrule
PBH $10^{12}$~kg & $1{,}2\times10^{11}$~K \\
Erde & $2{,}1\times10^{-2}$~K \\
Sonne & $6{,}2\times10^{-8}$~K \\
Sgr~A$^{*}$ & $1{,}5\times10^{-14}$~K \\
SMBH $10^{9}M_\odot$ & $6{,}2\times10^{-17}$~K \\
\bottomrule
\end{tabular}
\end{center}

Die Seifenblase aus Kap.~04 hat ein Thermometer, und das
Thermometer misst die Uhrenlandkarte.

\section{Die Familienleiter [S]}

Wo trifft die Leiter den Kosmos? $T_H = T_\text{GH}$ verlangt
$M = c^3/(4GH_0) = 4{,}66\times10^{52}$~kg -- und dessen
Horizont liegt bei

\[
  r_s = \frac{2GM}{c^2} = \frac{R_H}{2} \quad\text{(exakt } 0{,}500\text{)},
\]

zugleich exakt die \emph{halbe} kritische Hubble-Masse
$c^3/(2GH_0)$. Die Temperaturleiter läuft von Mikro-Löchern bis
zum Ganzen durch, ohne Bruch: Der Kosmos ist das größte
Familienmitglied.

\section{Was Hawking nicht leisten kann [K]}

Die ehrliche Reibung mit dem Zyklus: Verdampfungszeiten gegen
$\tau_c = 91{,}9$~Gyr:

\begin{center}
\begin{tabular}{lll}
\toprule
Objekt & $t_\text{evap}$ & $t_\text{evap}/\tau_c$ \\
\midrule
PBH $10^{12}$~kg & $10^{12{,}4}$~yr & $10^{1{,}5}$ \\
Sonne & $10^{67{,}3}$~yr & $10^{56{,}4}$ \\
SMBH $10^{9}M_\odot$ & $10^{94{,}3}$~yr & $10^{83{,}4}$ \\
\bottomrule
\end{tabular}
\end{center}

Nur primordiale Löcher unterhalb der Grenzmasse

\[
  M^{*} = \left(\frac{\tau_c\,\hbar c^4}{5120\,\pi G^2}\right)^{1/3}
  = 3{,}3\times10^{11}~\text{kg} \quad(\sim\text{Asteroid})
\]

schließen via Hawking innerhalb eines Zyklus. Ein
Sonnenmassen-Loch ist $56$~Größenordnungen zu langsam, ein SMBH
$83$. Die FFGFT-Korrektur aus Kap.~04,
$P = P_\text{std}(1 - \xi\ln(M/M_P))$, beträgt $0{,}6$--$1{,}4\,\%$
und ändert am Befund nichts.

\textbf{Konsequenz für D(ii):} Stellare und supermassive Löcher
-- die Träger des $10^{104}$-Entropiebudgets aus Kap.~D -- können
auf dem Zyklus \emph{nicht} via Hawking verschwinden. Die
Rückkehr-Hälfte muss sie auf einem anderen Weg zurückbauen.
Hawking-Strahlung passt also thermodynamisch perfekt ins Bild und
\emph{verschärft} zugleich den grobkörnigen Kern von D(ii): Sie
zeigt, wie hoffnungslos langsam der einzige bekannte Rückweg wäre.

\section{Informationsseite: zwei Sprachen, eine Aussage [S]}

Kap.~04 sagte immer schon: Die Strahlung korreliert mit der
fraktalen Membranstruktur -- nichts geht verloren
($S_\text{BH}(M_\odot) = 1{,}05\times10^{77}\,k_B$ an
Korrelationen). Kap.~D dieses Dokuments sagt: Feinkörnige Entropie
ist unter unitärer Evolution konstant -- Periodizität ist
mikroskopisch zulässig. Das ist dieselbe Aussage in zwei Sprachen:
Informationserhalt lokal (Membran) und global (Zyklus) sind
konsistent. Offen bleibt allein die \emph{grobkörnige}
Schließung -- und die war schon vorher der Kern.

% ============================================================
\chapter{Prüfflächen [S]}
% ============================================================

Photonen enden an der Wand ($0{,}992\,\pi$). Zwei Boten stammen
vom Antipoden selbst oder dahinter:

\begin{center}
\begin{tabular}{lll}
\toprule
Bote & Entkopplung & Position \\
\midrule
C$\nu$B (Neutrino-Hintergrund) & $z \sim 6\times10^{9}$ &
  $\theta \approx 1{,}01\,\pi$ \\
primordiale GW & nominell $z \sim 10^{25}$ &
  $\theta \approx 1{,}01\,\pi$ \\
\bottomrule
\end{tabular}
\end{center}

\textbf{Was die Wegkorrektur hier kostet:} Die Streuwand liegt
fraktal bei $0{,}979\,\pi$, also deutlich \emph{unter} dem
Halbzyklus -- sie schneidet sich nicht selbst. Das macht die
kompakte Lesart einerseits verträglicher mit den Nullresultaten
der Topologiesuchen (Schranke $(L^{*}/2)/D_\text{LSS}$ steigt von
$1{,}008$ auf $1{,}022$), nimmt ihr andererseits eine
Prüfsignatur: Ohne Selbstschnitt gibt es keine Matched Circles und
damit auch keinen $4{,}2'$-Versatz (Dok.~312). Der Preis der
besseren Zahlen ist eine verlorene Prüfmöglichkeit -- das ist
ehrlich zu buchen, nicht zu verschweigen.

Eine \emph{nicht-monotone} Signatur dort -- die Umkehr des
Massenlaufs jenseits des Antipoden -- bleibt damit die
\emph{einzige} verbliebene Prüffläche der Schließung. Praktisch außerhalb heutiger
Messbarkeit, aber benennbar; zusammen mit dem
$4{,}2'$-Dipolversatz (Dok.~312, Torus-Ruhesystem) sind das zwei
konkrete Signaturen der kompakten Lesart.

% ============================================================
\chapter{Die Antipoden-Bedingung als Vorhersage:
$\Omega_m^{*}$ und der $\kappa$-Kandidat}
% ============================================================

\section{Die Umkehrung}

Kap.~B hat die Lage der Streuwand als Befund gelesen und die
Koinzidenz $\eta_0 H_0/c \approx \pi$ ehrlich als [S] gebucht.
Die Bedingung lässt sich aber \emph{umkehren}: Wenn die kompakte
Lesart verlangt, dass der Partikelhorizont exakt am Antipoden
schließt,

\begin{equation}
  \eta_0 = \pi R_H
  \qquad\Longleftrightarrow\qquad
  \int_0^\infty \frac{dz}{E(z)} = \pi,
\end{equation}

dann ist $\Omega_m$ kein freier Parameter mehr. Die Bedingung
löst sich (flach, $\Omega_r = 9{,}3\times10^{-5}$) eindeutig zu:

Ohne fraktale Korrektur ergibt das
$\Omega_m^{*} = 0{,}325$ ($+1{,}4\,\sigma$ zu Planck). Mit ihr
lautet die Bedingung $\int dz/E = \pi/K_\text{frak} = 3{,}1841$
und liefert

\begin{equation}
  \boxed{\;\Omega_m^{*} = 0{,}3139\;}
  \qquad\text{gegen Planck } 0{,}315 \pm 0{,}007
  \;\;\Rightarrow\;\; -0{,}16\,\sigma.
\end{equation}

\noindent Die R67-Unsicherheit des Faktors ($n = 100\pm2$)
verschiebt das nur zwischen $0{,}3137$ und $0{,}3141$ -- die
Vorhersage ist robust gegen die einzige bezifferte Freiheit.

Der kosmische Sektor erhält damit erstmals eine \emph{Vorhersage
für einen $\Lambda$CDM-Parameter} statt einer Umdeutung --
falsifizierbar: Präzisere $\Omega_m$-Messungen (DESI, Euclid)
entscheiden, ob $0{,}325$ hält.

\textbf{Was die Vergleichsgröße ist:} Planck
$\Omega_m = 0{,}315 \pm 0{,}007$ ist kein Rohdatum, sondern
selbst der Output eines Sechs-Parameter-$\Lambda$CDM-Fits an
Photonenzählraten -- inklusive $\Lambda$CDM-Transferfunktionen
und Fiducial-Kosmologie in den Auswerteketten. Der
$1{,}4\,\sigma$-Vergleich ist also präzise gelesen:
\emph{Bedingung gegen Pipeline, nicht Bedingung gegen Natur.}
Das ist die dritte Vererbungsebene neben dem SI-Anker (Dok.~309)
und der $E(z)$-Form (oben), und sie wird wie diese deklariert
statt versteckt. Sie schneidet in beide Richtungen: Die
Jahrzehnte der $\Lambda$CDM-Abrundung erzeugen Robustheit
\emph{und} Pipeline-Lock-in -- und wo die kompakte Lesart die
Form erbt, erbt sie zugleich deren Fit-Erfolge für die geteilten
Teile. Deshalb liegen die eigenständigen Prüfpunkte dieses
Dokuments bewusst in der minimal modellbeladenen
Observablen-Klasse: Richtungen und Formen
($4{,}2'$-Dipolversatz, Nicht-Monotonie am Antipoden) statt
Fit-Parameter.

\textbf{Deklarierte Erbschaft:} Die Rechnung benutzt die
$\Lambda$CDM-\emph{Form} von $E(z)$ -- inklusive des
$(1{+}z)^3$-Terms mit vollem $\Omega_m$. Die Vorhersage ist also
präzise formuliert: \emph{Die Landkarte muss so geformt sein,
dass $\int dz/E = \pi$; in $\Lambda$CDM-Buchführung heißt das
$\Omega_m = 0{,}325$.} Was den Materie-Anteil dieser Buchführung
trägt, ist eine separate Frage (Abschn.~F.3). Status damit
[K$|$P20\,$\times$\,$E(z)$-Form] -- dieselbe
Vererbungsstruktur wie in Dok.~309, eine Ebene tiefer.

\section{Konsolidierung: Pflicht B und die Koinzidenz werden eine
Bedingung}

Aus $\Omega_m^{*}$ folgt $\Omega_\Lambda^{*} = 0{,}675$ und damit
ein struktureller Kandidat für den bisher offenen
Ordnung-1-Faktor der Abschlussskala (Dok.~312, Pflicht~B):

\begin{equation}
  \kappa^{*} = 3\,\Omega_\Lambda^{*} = 2{,}058
  \qquad\text{gegen gemessen } \kappa = 2{,}12
  \;\;(-2{,}9\,\%)
\end{equation}

\noindent (glatt gerechnet: $2{,}026$, also $-4{,}4\,\%$).

Dieselbe Konsolidierungslogik wie beim $10^{123}$-Problem
($\to$~P20): Zwei offene Stellen -- Pflicht~B und die
[S]-Koinzidenz aus Kap.~B -- werden auf \emph{eine} Bedingung
zurückgeführt. Die Gegenprobe zeigt, dass dies keine
Zirkelrechnung ist: $\kappa = 2{,}12$ würde $\Omega_m = 0{,}294$
verlangen, und dann wäre $\eta_0 H_0/c = 3{,}268 \ne \pi$.
Antipoden-Exaktheit und heutiger Planck-$\kappa$ stehen in einer
echten $4\,\%$-Spannung -- Prüfkontakt, kein Fit.

\section{Einordnungen und eine neue offene Stelle}

\textbf{$\Omega_\Lambda$ verliert den Substanzstatus:} Die
\enquote{Dunkle-Energie-Dichte} ist die FLRW-Übersetzung der
Abschlussskala; trägt die Antipoden-Bedingung, ist sie
vorhergesagt, nicht gefittet. \textbf{Hubble-Spannung:} Die
kompakte Lesart positioniert sich eindeutig beim CMB-verankerten
Wert ($H_0 = 66{,}8$ aus der $\xi^{10}$-Kette gegen Planck
$67{,}4$); der lokale SH0ES-Überschuss müsste ein
Landkarten-Effekt der Kalibration sein [S].

\textbf{Dunkle Materie: zwei Regime, eines offen.} Von Dunkler
\emph{Materie} war bisher bewusst nicht die Rede -- und die Lücke
gehört benannt. \emph{Galaktisch} (Rotationskurven, RAR) braucht
die FFGFT keine Teilchen-DM: Dort trägt die Übergangsskala $a_0$
(Dok.~308); davon sagt dieses Dokument nichts, und das ist
korrekt. \emph{Auf Hintergrundebene} aber erbt die
$z\to\chi$-Landkarte den CDM-Anteil stillschweigend mit: In
$\Lambda$CDM-Buchführung besteht $\Omega_m^{*} = 0{,}325$ aus
Baryonen ($\approx 0{,}049$) plus kalter Dunkler Materie
($\approx 0{,}276$). Was diesen Anteil in der kompakten Lesart
trägt -- tatsächliche Materie oder eine Eigenschaft des
Massenlauf-Profils $m(\theta)$, die in FLRW-Übersetzung wie
Materie bucht -- ist ungeklärt und wird als \textbf{zweite neue
offene Stelle} gebucht. Sie ist von der galaktischen Frage
unabhängig und darf nicht mit ihr verwechselt werden.

\textbf{Präzisierte offene Stelle: $T_0$ ist zweifach gebunden,
aber nicht hergeleitet.} Was setzt die Monopol-Temperatur
$T_0 = 2{,}725$~K? Der Korpus enthält dazu zwei Bindungen auf sehr
verschiedenem Niveau. \emph{(a)~Relational (Dok.~279, kanonisch):}
$T_0$ definiert die Casimir-Länge
$L_\xi = (15\xi/\pi^2)^{1/4}\,\hbar c/(k_BT_0) = 100{,}24~\mu$m
mit dem $\xi$-Verhältnis
$\rho_\text{Casimir}/\rho_\text{CMB} = \pi^2/(240\xi) = 308$;
$T_0$ geht dort als \emph{gemessener Input} ein, und Dok.~279
warnt selbst, dass $m_ec^2/(k_BT_0)$ ein separater,
nicht-geometrischer Faktor ist. \emph{(b)~Vorwärts-Behauptung
(Dok.~061, alte Schicht):} $k_BT_0 = (16/9)\,\xi$~eV trifft auf
$+0{,}92\,\%$ -- ist aber nach heutigem Maßstab nicht bookbar:
Der eV ist ein verborgener SI-Anker (der Faktor
$1/510\,999$ zur Kammerton-Form ist nicht strukturell,
Verstoß gegen die Dok.-310-Disziplin), der Faktor $16/9$ ist
post-hoc, und die dortige Präzisionsangabe
(\enquote{$0{,}0002\,\%$}) ist falsch (real $0{,}92\,\%$) --
Korrekturregister-Kandidat. \emph{Fazit:} Dimensionslos ist
$k_B T_0/(m_e c^2) = 4{,}6\times10^{-10}$ keine ganze
$\xi$-Potenz ($\log_\xi = 2{,}41$); nach R61-Maßstab wird nicht
numerologisiert. $T_0$ bleibt \emph{vorwärts unhergeleitet} und
als offene Frage des kosmischen Sektors gebucht -- relational
eingebunden, aber ohne Träger.

% ============================================================
\chapter{Statusübersicht}
% ============================================================

\begin{center}
\begin{tabular}{lp{6.8cm}l}
\toprule
Aussage & Inhalt & Status \\
\midrule
Antipoden-Befund & Horizont am Antipoden: fraktal $0{,}14\,\%$
  (glatt $1{,}2\,\%$); Streuwand $0{,}979\,\pi$ & [K$|$P20] \\
Wegkorrektur & $K_\text{frak} = 1-100\xi$ greift bei Weg/Struktur,
  kürzt bei Weg/Weg (Test bestanden) & [K]/[S] \\
Koinzidenz-Einordnung & $\eta_0 H_0/c \approx \pi$: Strukturaussage
  statt Zufall erst nach Vorwärtsherleitung & [S] \\
Glatte Schließung & Endpunktsteigung $\sqrt{\Omega_r}/K
  \approx 0{,}0098$; Schließung auf $1\,\%$ erfüllt & [K] \\
Rückkehr-Hälfte & $=$ photonisch unbeobachtbare Hälfte;
  kein Konflikt mit Photon-Daten & [K] \\
Periodizität & Randbedingung (Poincaré-Diskrepanz
  $10^{104}$~dex im Exponenten), nicht Dynamik & [K] \\
Entropie & feinkörnig zulässig; grobkörnige Schließung offen --
  präzisierter Kern von D(ii) & offen \\
Hawking & lokaler Fall derselben KMS-Regel (GH aus Dok.~312);
  Leiter trifft Kosmos bei $r_s = R_H/2$ & [K]/[S] \\
Verdampfung & $M^{*} = 3{,}3\times10^{11}$~kg; $M_\odot$: $56$~dex
  zu langsam -- verschärft D(ii) & [K] \\
Prüfflächen & C$\nu$B/GW vom Antipoden; Nicht-Monotonie als
  Signatur & [S] \\
$\Omega_m^{*}$ & fraktal $0{,}3139$ ($-0{,}16\,\sigma$), glatt
  $0{,}325$ ($+1{,}4\,\sigma$); $E(z)$-Form geerbt
  & [K$|$P20$\times$$E$] \\
$\kappa^{*}$ & fraktal $2{,}058$ ($-2{,}9\,\%$), glatt $2{,}026$:
  Kandidat für Pflicht~B & [S] \\
Matched Circles & Streuwand ohne Selbstschnitt: Signatur aus
  Dok.~312 entfällt & Verlust \\
DM (Hintergrund) & $(1{+}z)^3$-Anteil der Landkarte: Träger
  ungeklärt (galaktisch: $a_0$, Dok.~308, davon getrennt) & offen \\
$T_0$ & relational gebunden (Dok.~279: $L_\xi$, $\pi^2/240\xi$);
  061er-Relation eV-verankert, nicht bookbar; vorwärts offen
  & offen \\
\bottomrule
\end{tabular}
\end{center}

\medskip
\noindent\textbf{Bilanz:} \enquote{Kein Urknall} heißt nach diesem
Dokument präzise: Der \enquote{Anfang} ist der Antipode des
Zeitzyklus -- ein Ort, kein Ereignis; die beobachtete Geschichte
ist die eine Hälfte einer auf $1\,\%$ glatt schließbaren Landkarte
(mit fraktaler Wegkorrektur trifft der Partikelhorizont den
Antipoden auf $0{,}14\,\%$);
und die verbleibende Last liegt nicht bei den Feldern, sondern bei
der grobkörnigen Entropie. D(ii) ist damit nicht erledigt, sondern
auf seinen harten Kern reduziert.

\medskip
\noindent\textbf{Registrierung:} Aufnahme in Dok.~190 erfolgt
separat.

\end{document}
